\documentclass[runningheads]{llncs}
\usepackage{graphicx}

\begin{document}

\title{A Reactive Swarm Controller for the ARGoS Tunnelling Task}

\author{Louis Cours}

\authorrunning{L. Cours}

\institute{Swarm Intelligence, INFO-H-414}

\maketitle

\begin{abstract}
This report presents a reactive controller for the ARGoS tunnelling project.
The implementation combines local repulsion and attraction vectors, a
phase-based state machine, and a small set of arena cues to coordinate a swarm
of foot-bots without centralized control. The controller uses range-and-bearing
communication to maintain spacing, proximity and light sensing to bias motion,
and motor-ground readings to detect the black target zone. The current version
implements three main phases: initial orientation, grouping, and tunnel
progression, followed by a black-zone behavior that keeps the robots inside the
target region once they enter it. The report also documents the project
requirements, the design choices made in the controller, and the evaluation
protocol to be used for scaling experiments. Experimental figures and summary
tables are left as explicit placeholders where results still need to be added.
\keywords{swarm robotics \and ARGoS \and Lua \and tunnelling \and foot-bot}
\end{abstract}

\section{Introduction}
The project asks for a cooperative collective behavior able to solve the
tunnelling task in ARGoS: the swarm must move through an obstacle-filled arena,
reach the target zone on the right, and avoid bringing obstacles into the
target. The score is defined as the number of robots in the target minus the
number of obstacles in the target at the end of the run. The assignment also
requires repeated runs with different random seeds, performance summaries in the
form of boxplots, and a scalability analysis across swarm sizes.

The implementation discussed in this report is written in Lua for the foot-bot
platform and follows the local sensing and actuation model described in the
course notes for Lua and the foot-bot API~\cite{lua,footbot}. The design is
intentionally reactive: each robot reacts to nearby peers, light, floor color,
and obstacles rather than relying on a global map or centralized coordination.

\section{Problem Definition and Requirements}
The arena is a rectangular environment with a white starting area on the left,
an obstacle field in the middle, and a black target area on the right. The
project brief specifies a 1{,}000 second run, corresponding to 10{,}000 ARGoS
steps, although a shorter runtime may be used if justified in the report.

The available sensors and actuators include proximity, light, motor-ground,
range-and-bearing, colored-blob omnidirectional camera, wheels, LEDs, gripper,
and turret. The brief emphasizes cooperative behavior, use of swarm robotics
principles, repeated runs over 10 seeds, and a scaling study over several swarm
sizes. In addition to the final score, the project data file records the number
of robots and obstacles in the target at each step.

\subsection{Evaluation Criteria}
The main quantities to report are:
\begin{itemize}
	\item final score as a function of swarm size,
	\item number of robots in the target over time,
	\item number of obstacles in the target over time,
	\item variability across the 10 repetitions for each configuration.
\end{itemize}

\section{Controller Overview}
The controller is organized as a state machine with four phases: orientation,
grouping, tunnel progression, and black-zone behavior. The transition logic is
driven by local conditions such as the number of neighbors at a target spacing,
floor color, and whether the robot has already entered the black target zone.

At a high level, the controller combines three vector sources:
\begin{itemize}
	\item a Lennard--Jones interaction over range-and-bearing messages,
	\item a light-seeking vector from the ambient light sensor,
	\item obstacle repulsion and tangential corrections from proximity readings.
\end{itemize}

The implementation also uses robot LEDs to expose the current phase, which is
useful during debugging and qualitative analysis. Per-robot logs are written to
the \texttt{logs/} directory.

\subsection{Orientation Phase}
In the initial phase, robots bias motion toward the light while maintaining some
distance from nearby robots and obstacles. The controller uses the light vector,
the Lennard--Jones interaction, and a small obstacle tangential correction to
avoid overcrowding at the start. After a fixed number of steps, the robot moves
to the grouping phase.

\subsection{Grouping Phase}
The grouping phase uses the range-and-bearing system to enforce a loose swarm
structure. The code computes a Lennard--Jones potential at the current peer
distance and counts neighbors in a target interval around 80 cm. Once enough
neighbors are found consistently for a sufficient number of steps, the robot
switches to the tunnel phase.

The controller also computes a leader vector from robots that already advanced
into later phases. This creates a simple local bootstrap effect: robots that are
already closer to the target influence nearby robots to follow them.

\subsection{Tunnel Phase}
The tunnel phase reduces the target spacing to 60 cm and strengthens the
attraction toward robots already in the tunnel. The phase uses a combination of
Lennard--Jones interaction, light attraction, and leader following. If an
obstacle is directly ahead, the controller adds a sideways correction so the
robot can move around the obstacle field instead of pushing straight through it.

This phase is the main navigation behavior of the current implementation. It is
designed to help the swarm advance collectively through the clutter while still
preserving a coherent formation.

\subsection{Black-Zone Phase}
When enough floor sensors detect black for several consecutive steps, the robot
enters the black-zone state. In this state, the controller keeps the robot in
the target region and uses a randomized heading segment to avoid getting stuck
on the boundary. If the robot sees the green beacon of another robot already in
the black zone, it will bias motion toward that robot as a simple local cue.

The motor-ground readings are central here: they determine when the robot has
reached the target and when it should leave the boundary-following motion.

\section{Implementation Details}
\subsection{Sensors and Local Rules}
The controller uses the following signals:
\begin{itemize}
	\item \textbf{Proximity sensors}: obstacle repulsion, obstacle tangential motion, and wall/robot blocking.
	\item \textbf{Light sensors}: orientation toward the ambient cue.
	\item \textbf{Motor-ground sensors}: detection of the black target zone.
	\item \textbf{Range-and-bearing}: local communication and Lennard--Jones spacing.
	\item \textbf{Colored-blob camera}: detection of a green beacon emitted by robots already in the black zone.
	\item \textbf{LEDs}: visual phase feedback during simulation.
\end{itemize}

The wheel command is computed from the target vector angle with a simple
proportional steering law. The result is clamped to the allowed wheel speed.
This keeps the behavior easy to reason about and consistent with the reactive
style encouraged in the project brief.

\subsection{State Transitions}
The state machine is driven by local counts and thresholds rather than by global
knowledge. The most relevant conditions are:
\begin{itemize}
	\item a minimum number of neighbors in the target spacing interval,
	\item a minimum duration during which that grouping condition must hold,
	\item a consecutive black-floor counter before entering the black zone,
	\item a beacon-based check to bias motion once inside the target region.
\end{itemize}

\subsection{Current Limitation}
The controller contains an obstacle-grabbing sequence using the gripper and
turret, but that path is currently disabled in the implementation because the
corresponding handler returns immediately. As a result, the current version
focuses on navigation, grouping, and black-zone stabilization rather than on
physical obstacle transport. This should be stated explicitly in the final
report, since the project brief expects the gripper and turret to be used for a
full tunnelling solution.

\section{Experimental Protocol}
The project brief requests 10 runs per configuration with different random
seeds. The report should present boxplots of the final score for each swarm
size, along with time-series plots for the number of robots in the target and
the number of obstacles in the target.

\subsection{Planned Settings}
\begin{table}[htbp]
\centering
\caption{Planned experimental settings. Replace placeholders with the final tested configurations.}
\label{tab:settings}
\begin{tabular}{|l|l|l|}
\hline
Configuration & Value & Notes \\
\hline
Run duration & 10{,}000 steps & Default setting \\
\hline
Seeds per setting & 10 & Required \\
\hline
Swarm sizes & \texttt{<insert sizes>} & Prefer higher resolution \\
\hline
Baseline & \texttt{<insert baseline if any>} & Optional comparison \\
\hline
\end{tabular}
\end{table}

\subsection{Data to Report}
\begin{table}[htbp]
\centering
\caption{Summary placeholders for experiment outputs. Replace with computed statistics once the runs are available.}
\label{tab:summary}
\begin{tabular}{|l|l|l|l|}
\hline
Swarm size & Mean score & Std. dev. & Median score \\
\hline
\texttt{<n1>} & \texttt{<value>} & \texttt{<value>} & \texttt{<value>} \\
\hline
\texttt{<n2>} & \texttt{<value>} & \texttt{<value>} & \texttt{<value>} \\
\hline
\texttt{<n3>} & \texttt{<value>} & \texttt{<value>} & \texttt{<value>} \\
\hline
\end{tabular}
\end{table}

\section{Results}
The following figures are placeholders for the final analysis. They should be
replaced by the actual plots generated from the 10 repeated runs per swarm size.

\begin{figure}[htbp]
\centering
\fbox{\parbox[c][4.2cm][c]{0.92\textwidth}{\centering
Placeholder for the boxplot of the final score versus swarm size.}}
\caption{Final score as a function of swarm size. Each box should summarize 10 repetitions.}
\label{fig:score-boxplot}
\end{figure}

\begin{figure}[htbp]
\centering
\fbox{\parbox[c][4.2cm][c]{0.92\textwidth}{\centering
Placeholder for the time evolution of the number of robots in the target, with variance across runs.}}
\caption{Robots in the target over time for representative swarm sizes.}
\label{fig:robots-over-time}
\end{figure}

\begin{figure}[htbp]
\centering
\fbox{\parbox[c][4.2cm][c]{0.92\textwidth}{\centering
Placeholder for the time evolution of the number of obstacles in the target, with variance across runs.}}
\caption{Obstacles in the target over time for representative swarm sizes.}
\label{fig:obstacles-over-time}
\end{figure}

\subsection{Discussion Points for the Final Draft}
Once the data are inserted, the discussion should answer the following:
\begin{itemize}
	\item How does performance scale with the number of robots?
	\item Is the performance monotonic, or do some sizes perform worse than nearby sizes?
	\item Do some swarm sizes exhibit higher variance than others?
	\item Are the observed trends consistent with the controller's local rules and phase transitions?
\end{itemize}

\section{Conclusion}
The current controller implements a clear reactive structure that is aligned
with the swarm robotics principles emphasized in the course: local sensing,
distributed coordination, and simple rules that produce collective motion.
The report draft is complete enough to receive the final experimental plots and
statistics, but the obstacle-transport capability still needs to be finished or
explicitly presented as an open limitation if no further code changes are made.

\begin{thebibliography}{8}
\bibitem{lua}
H414 -- Swarm Intelligence: Programming the Robots with Lua. Course reference document.

\bibitem{footbot}
H414 -- Swarm Intelligence: Sensor and actuator reference for the foot-bot. Course reference document.

\bibitem{projectbrief}
H414 -- Swarm Intelligence: Evaluation and Final Project of INFO-H-414 Swarm Intelligence, Second Session.

\bibitem{llncs}
Springer: Lecture Notes in Computer Science style and sample paper documentation.
\end{thebibliography}
\end{document}
